\chapter{Introduction}
\label{ch:intro}

\section{Motivation}

Autonomous aerial systems are moving from single, monolithic vehicles toward
\emph{clusters}: coordinated groups of many UAVs that share sensing, computation
and objectives. A cluster survives the loss of individual members, covers more
ground per unit time, localises targets by triangulation that no single vehicle
can achieve, and reshapes itself on demand into search lines, perimeter rings or
convoy escorts. Realising these behaviours in the field forces three problems
together that are usually studied in isolation.

First, the vehicles are \emph{dynamically heterogeneous}. A multirotor
hexacopter hovers, translates in any direction and reorients almost
instantaneously, but is energetically expensive and slow; a fixed-wing aircraft
is efficient and fast but cannot stop, must maintain airspeed above stall and
turns only by banking. A framework that is to command a mixed fleet must model
both to the same fidelity and expose a common control interface. Second, the
vehicles operate under \emph{hard physical constraints}---thrust and actuator
limits, minimum airspeeds, current and voltage envelopes, no-fly volumes and,
above all, mutual collision avoidance---that are safety-critical and active
exactly when the mission is most demanding. Third, coordination must be
\emph{distributed and real-time}: bandwidth is limited, a central node is a
single point of failure, and every vehicle must close its own estimation and
control loops at rates measured in hundreds of hertz on an embedded
companion computer.

This report addresses the three problems jointly. It develops complete,
first-principles dynamic models of a quaternion-attitude hexacopter and of a
six-degree-of-freedom fixed-wing aircraft; it places at the core of each agent a
constrained real-time nonlinear model-predictive \emph{controller} and a
constrained real-time moving-horizon \emph{estimator}; and it builds around that
core a distributed estimation-and-control architecture that couples the agents
over a communication graph into safe, coordinated formations.

\section{Problem statement}

Consider a cluster of $N_a$ UAVs indexed by $i\in\mathcal V=\{1,\dots,N_a\}$,
each with nonlinear continuous-time dynamics
\begin{equation}
\dot{\state}_i = \vv{f}_i(\state_i,\ctrl_i), \qquad
\vv{y}_i = \vv{h}_i(\state_i) + \vv{n}_i,
\label{eq:agent-generic}
\end{equation}
where $\state_i$ is the vehicle state, $\ctrl_i$ its control input, $\vv{y}_i$ a
vector of noisy onboard and inter-agent measurements, and $\vv{n}_i$ measurement
noise. The vehicle class of agent $i$---hexacopter or fixed-wing---fixes the
form of $\vv{f}_i$ and $\vv{h}_i$, derived in
\cref{ch:hexacopter,ch:fixedwing}. The agents share a time-varying communication
graph $\Graph=(\mathcal V,\mathcal E)$; agent $i$ may exchange messages only with
its neighbours $\Neigh_i=\{j:(i,j)\in\mathcal E\}$. Each agent must, using only
local measurements and neighbour messages,
\begin{enumerate}[nosep,leftmargin=1.4em]
  \item \emph{estimate} its own state $\state_i$ (and unmeasured disturbances)
        in real time, remaining accurate even when a subset of its
        absolute-positioning measurements is lost;
  \item \emph{control} its motion to track a share of a cluster-level objective
        (a formation, a trajectory, a coverage pattern) while respecting its own
        actuator and state constraints; and
  \item \emph{coordinate} with neighbours so that the cluster as a whole is
        collision-free, connected and consistent.
\end{enumerate}
The design target throughout is a per-agent computation and communication budget
small enough to run on an embedded companion computer at the vehicle's control
rate.

\section{Approach and contributions}

The architecture is layered, and the layers map onto the natural language
boundaries of an implementation (\cref{ch:architecture}):
\begin{itemize}[nosep,leftmargin=1.4em]
  \item a \emph{numerical core}, in C, providing the two constrained real-time
        solvers---the MPC controller and the MHE estimator---as the per-agent
        engine (\cref{ch:peragent});
  \item a \emph{coordination layer}, in C++, implementing the distributed
        estimation (\cref{ch:dmhe}) and distributed control
        (\cref{ch:dmpc,ch:formation}) around that core, together with
        neighbour discovery, consensus, message passing and fault handling; and
  \item a \emph{research layer}, in Python/MATLAB, for simulation, tuning and
        analysis (\cref{ch:simulation}).
\end{itemize}
The specific contributions are the following.
\begin{enumerate}[leftmargin=1.6em]
  \item \textbf{A complete quaternion hexacopter model
        (\cref{ch:hexacopter})}, derived from first principles: Hamilton
        quaternion kinematics, Newton--Euler translational and rotational
        dynamics, the rotor thrust/torque model, and the six-rotor control
        allocation with its rank structure, invertibility and fault tolerance,
        followed by discretisation, analytic Jacobians and a differential-%
        flatness map. All relations are validated numerically.
  \item \textbf{A complete fixed-wing model (\cref{ch:fixedwing})} in the
        Beard--McLain six-degree-of-freedom formulation, including the wind
        triangle, the sigmoid-blended stall aerodynamics, the full force and
        moment coefficient expansions, the product-of-inertia rotational
        dynamics, and a numerically solved trim, on the fully parameterised
        Aerosonde airframe.
  \item \textbf{A generic account of the per-agent constrained real-time
        estimator and controller (\cref{ch:peragent})} as a matched
        estimation--control pair that imposes general nonlinear state and input
        constraints on the full problem and executes one real-time iteration per
        sample, exposed to the coordination layer through a clean C application
        programming interface. The solvers' internal algorithms are separate
        unpublished work and are treated here as a fixed, validated capability
        specified by that interface (see the scope note in \cref{ch:peragent}).
  \item \textbf{A distributed moving-horizon estimator (\cref{ch:dmhe})} in
        which agents fuse relative inter-agent measurements to their
        well-localised neighbours, with a covariance-aware fusion rule that
        keeps GNSS-denied agents accurate through absolute-positioning outages.
  \item \textbf{A distributed model-predictive controller
        (\cref{ch:dmpc,ch:formation})} that couples the agents through formation
        objectives and control-barrier-function collision-avoidance constraints,
        coordinated by consensus and alternating-direction methods, with the
        cooperative-DMPC stability and recursive-feasibility properties stated
        and attributed.
  \item \textbf{A reference implementation and simulation study
        (\cref{ch:architecture,ch:simulation})} on the full nonlinear plants,
        exercising formation assembly, reconfiguration, collision avoidance and
        GNSS-denied localisation for both vehicle classes.
\end{enumerate}

\section{Report organisation}

\Cref{ch:preliminaries} fixes the frames, rotation representations, rigid-body
mechanics and graph-theoretic tools used throughout.
\Cref{ch:hexacopter,ch:fixedwing} derive the two vehicle models.
\Cref{ch:peragent} describes the per-agent estimator and controller.
\Cref{ch:dmhe,ch:dmpc,ch:formation} build the distributed estimation, control
and formation layers. \Cref{ch:architecture} sets out the software architecture
and \cref{ch:simulation} the simulation study. \Cref{ch:conclusion} concludes.
Appendices collect quaternion identities, the full hexacopter Jacobians and the
complete vehicle parameter tables.

\keybox{\textbf{A note on the per-agent solvers.} Throughout this report the two
onboard solvers are referred to generically as ``the MPC controller'' and ``the
MHE estimator''. Their internal algorithms are treated as a fixed, validated
capability accessed through the interface of \cref{ch:peragent}; the distributed
framework depends only on that interface and not on the solvers' internal
construction.}
